\hypertarget{question}{}
\hypertarget{the-research-question}{%
\chapter{The research question}\label{the-research-question}}

The study examines 25 portfolios formed from the intersections of five
size groups and five book-to-market groups. It asks whether their
returns contain persistent performance unexplained by a standard factor
benchmark, how uncertain any resulting ranking is, and whether a ranking
estimated from the past predicts genuinely future benchmark-adjusted
returns. This chapter states those questions precisely, fixes what would
count as an answer to each, and explains why the exercise is far harder
than running twenty-five regressions and reading off the largest
intercept. The tools it invokes are all derived in Chapter~\ref{preliminaries};
here they are assembled into a single research design.

\section{Formalising the estimand and the three questions}\label{rq-estimand}

Let $i=1,\dots,N$ index the $N=25$ portfolios and $t=1,\dots,T$ the
$T=1{,}193$ months. For each portfolio the benchmark regression of
Chapter~\ref{preliminaries} is
\begin{equation}\label{eq:rq-model}
y_{it}=\alpha_i+\beta_i^{\top}f_t+\varepsilon_{it},\qquad
y_{it}=r_{it}-r_{ft},
\end{equation}
with $f_t$ the FF3 factor vector. The population object of interest is the
$N$-vector of intercepts $\alpha=(\alpha_1,\dots,\alpha_N)^{\top}$ and,
derived from it, the ranking that orders the portfolios from highest to
lowest. Everything the study reports is a functional of $\alpha$ and of
its estimation uncertainty. The research question decomposes into three
sharply different statistical problems.

\begin{longtable}[]{@{}p{0.06\textwidth}p{0.30\textwidth}p{0.28\textwidth}p{0.28\textwidth}@{}}
\toprule\noalign{}
& Question & Formal object & What counts as an answer \\
\midrule\noalign{}
\endhead
\bottomrule\noalign{}
\endlastfoot
Q1 & Is there benchmark-relative performance at all? & The vector
$\alpha$ and the joint null $H_0:\alpha=0$ & A dependence-robust rejection
of $H_0$ and per-portfolio evidence surviving multiplicity control \\
Q2 & How well-determined is the ranking? & The rank vector
$R=\operatorname{rank}(\hat\alpha^{\mathrm{post}})$ as a random object &
An honest sampling/resampling distribution of each rank, not a point list \\
Q3 & Does a past ranking predict future benchmark-adjusted returns, and
does re-estimating it help? & The forward association between a frozen
ranking and later alpha, and its \emph{increment} over a fixed ordering &
A strictly leakage-free, dependence-robust test, and a decomposition
separating dynamic from static content \\
\end{longtable}

Q1 is a question of \emph{inference} (Chapters~6--7); Q2 is a question of
\emph{estimation uncertainty} (Chapters~8--9); Q3 is a question of
\emph{prediction} (Chapters~11--12). They are logically independent: one
can reject $H_0:\alpha=0$ (Q1) yet find the ranking almost undetermined
(Q2) and its dynamic content null (Q3), which is in fact what happens.

\section{Why ``performance'' is model-relative}\label{rq-model-relative}

The estimand $\alpha_i$ is defined only \emph{relative to the chosen
factors}. Writing the population projection of $y_i$ on $(1,f)$, the
intercept is
\begin{equation}
\alpha_i=\mathbb{E}[y_{it}]-\beta_i^{\top}\mathbb{E}[f_t],
\end{equation}
so $\alpha_i$ is the average excess return left after removing exposure to
\emph{those} factors and nothing more. A non-zero $\alpha_i$ can arise from
an omitted priced risk, a friction, a nonlinear payoff, a change in factor
loadings over time, or genuine abnormal performance --- the regression
cannot tell these apart. On the 25 test assets the interpretation is even
narrower: these are the very portfolios from which the size and value
factors are constructed, so their $\alpha$ is a \emph{pricing residual} of
the model against its own building blocks, not a measure of skill. The
course therefore never writes ``skill'' for $\alpha$; it writes
``benchmark-relative'' and treats benchmark choice itself as a variable
(the FF3-versus-FF5 sensitivity of Chapter~13).

\section{Why this is harder than running 25 regressions}\label{rq-difficulties}

A naive analysis would run \eqref{eq:rq-model} twenty-five times, take the
largest $\hat\alpha_i$, and declare a winner. Seven features of the problem
make that invalid; each is quantified below and handled by a specific tool
from Chapter~\ref{preliminaries}.

\subsection{1. The alpha estimates are correlated}

The portfolios share calendar dates, common market shocks and overlapping
construction rules, so their residuals $\varepsilon_{it}$ are correlated
across $i$ within a month. By the influence-weight identity
\eqref{eq:influence}, $\hat\alpha_i-\alpha_i=\sum_t a_t\varepsilon_{it}$
with weights common to all portfolios, so the alpha errors inherit that
cross-sectional dependence and the estimator's covariance is a
\emph{full} matrix $V$, not a diagonal one. This matters concretely: the
precision of a comparison between two portfolios is
\begin{equation}
\operatorname{Var}(\hat\alpha_i-\hat\alpha_j)=V_{ii}+V_{jj}-2V_{ij},
\end{equation}
which positive covariance $V_{ij}$ can make far smaller --- or, if
negative, larger --- than the sum of the marginal variances. Twenty-five
separate standard errors discard every $V_{ij}$ and therefore misstate
both joint tests and rankings. The remedy is the joint HAC covariance of
\S\ref{prelim-hac} and Chapter~6.

\subsection{2. Selecting the largest estimate is a winner's curse}

Even if every true $\alpha_i$ were equal, the largest \emph{observed}
estimate would exceed the truth, because extreme draws combine signal with
favourable noise. In the simplest case $\hat\alpha_i=\alpha_i+e_i$ with
$e_i\sim N(0,\sigma^2)$ independent, the expected maximum of the noise
alone is
\begin{equation}
\mathbb{E}\Bigl[\max_{i\le N} e_i\Bigr]\ \approx\ \sigma\sqrt{2\ln N},
\end{equation}
so the top-ranked estimate is biased upward by roughly
$\sigma\sqrt{2\ln 25}\approx 2.54\,\sigma$ in this idealisation. Reporting
the raw maximum therefore overstates the leader and, symmetrically,
understates the laggard. The correction is partial pooling: the
empirical-Bayes shrinkage of \S\ref{prelim-eb} pulls noisy extremes toward
the common mean by the weight $\tau^2/(\tau^2+v)$ of
\eqref{eq:scalar-post}, and the James--Stein dominance result guarantees
this reduces risk in dimension $N\ge 3$.

\subsection{3. Testing twenty-five hypotheses inflates false discoveries}

Screening $N=25$ nulls at level $\alpha=0.05$ each makes at least one false
rejection likely even when all nulls hold: for independent tests
\begin{equation}
\Pr(\text{at least one false positive})=1-(1-0.05)^{25}\approx 0.72,
\end{equation}
and the expected number of false positives is $N\alpha=1.25$. A ``5\%
significant'' alpha found this way is uninformative on its own. The study
controls the false-discovery rate with the Benjamini--Hochberg procedure of
\S\ref{prelim-fdr}, reporting monotone $q$-values \eqref{eq:bh}, and keeps
the raw-alpha family separate from the frontier-test family because they
test different nulls.

\subsection{4. Ranks are discontinuous}

The rank of portfolio $i$,
\begin{equation}
R_i=1+\sum_{j\ne i}\mathbf 1\{\hat s_j>\hat s_i\},
\end{equation}
is a step function of the scores $\hat s$: it is flat almost everywhere and
jumps by an integer when two scores cross. Its derivative is zero wherever
it is defined, so the delta method \eqref{eq:delta} cannot supply a
standard error, and an arbitrarily small change in a score can move a rank
by several positions. Rank uncertainty must therefore be obtained by
resampling the \emph{entire} estimator, which is what the common-date
circular block bootstrap of \S\ref{prelim-boot} and Chapter~9 does; the
result is a rank interval, not a point.

\subsection{5. Overlapping windows manufacture persistence}

The rolling analysis uses windows of length $w=120$ months advanced by a
step $s=12$. Two windows $k$ steps apart share $\max(0,w-ks)$ months, an
overlap fraction
\begin{equation}
\frac{w-ks}{w}=1-\frac{ks}{w},\qquad ks<w,
\end{equation}
so \emph{adjacent} windows ($k=1$) share $108/120=90\%$ of their data. A
high correlation between adjacent rolling estimates is then largely
mechanical: the two estimates are computed from almost the same months.
Structural persistence can only be read from \emph{non-overlapping}
comparisons, which require $ks\ge w$, i.e.\ $k\ge w/s=10$ steps $=120$
months apart. Chapter~10 gives structural weight only to those disjoint
comparisons, where persistence turns out to be weak.

\subsection{6. A forward relation need not be dynamic skill}

This is the subtle point on which the study turns. Suppose each cell has a
\emph{fixed} benchmark-relative level $c_i$ (small-growth structurally
poor, small-value structurally strong) so that both the training score and
the future alpha are noisy readings of the same $c_i$:
\begin{equation}
\hat s_{it}=c_i+u_{it},\qquad A^{+}_{i,t}=c_i+w_{i,t}.
\end{equation}
Then the cross-sectional correlation $\operatorname{Corr}_i(\hat s_{it},
A^{+}_{i,t})$ is positive purely through the shared fixed effect $c_i$,
with \emph{no time-varying information} in the ranking at all. A positive
forward correlation is thus consistent with a merely \emph{stable}
ordering. Distinguishing this from genuine dynamic forecasting requires
comparing the time-varying rolling ranking against a fixed ordering on the
same future data --- the incremental test of \S\ref{prelim-forecast} and
Chapter~12, which is a Diebold--Mariano test of equal predictive accuracy
\eqref{eq:dm}.

\subsection{7. Alpha depends on the benchmark}

Because $\alpha_i=\mathbb{E}[y_i]-\beta_i^{\top}\mathbb{E}[f]$ depends on
which factors $f$ are included, enlarging the model changes both the
loadings and the spanned space and hence the intercept. A portfolio that
loads on a priced dimension absent from FF3 will show a spurious FF3 alpha
that shrinks or reverses under FF5. The study measures this directly: on a
common sample the FF3 and FF5 rankings correlate but a single portfolio
moves eleven rank positions (Chapter~13). Benchmark dependence is a
property of the estimand, not a nuisance to be tuned away.

\section{The central logic: stable ordering versus dynamic skill}\label{rq-logic}

The seven difficulties combine into one argument, which the reader should
be able to reconstruct and, above all, to stop at the correct conclusion.

\begin{enumerate}
\tightlist
\item
  Historical benchmark-relative alpha varies across the cross-section
  (description, Chapter~5).
\item
  A ranking formed from past alpha correlates with future
  benchmark-adjusted returns (forward validation, Chapter~11).
\item
  But by difficulty~6 this is also what a \emph{fixed} ordering would
  produce, so the rolling ranking is compared against static and
  expanding-window orderings on identical future data (incremental test,
  Chapter~12).
\item
  The rolling ranking adds no measurable increment over the fixed
  ordering: the Diebold--Mariano differential is statistically
  indistinguishable from zero.
\item
  Conclusion: the predictable component is a \emph{stable} cross-sectional
  ordering of persistent characteristic mispricing, \emph{not} dynamic
  timing skill.
\end{enumerate}

Formally, the study pits two hypotheses against each other. Let $\rho^{\mathrm{roll}}_t$
and $\rho^{\mathrm{fix}}_t$ be the window-$t$ predictive scores of the
rolling and fixed orderings and $d_t=\rho^{\mathrm{roll}}_t-\rho^{\mathrm{fix}}_t$.
\begin{equation}
H_{\text{stable}}:\ \mathbb{E}(d_t)\le 0
\qquad\text{versus}\qquad
H_{\text{dynamic}}:\ \mathbb{E}(d_t)>0 .
\end{equation}
The forward test of Q3 establishes only that
$\mathbb{E}(\rho^{\mathrm{roll}}_t)>0$; it is the incremental test that
adjudicates $H_{\text{stable}}$ against $H_{\text{dynamic}}$, and it fails
to reject $H_{\text{stable}}$. A student must be able to distinguish the
weak, supported claim ``a past ranking correlates with future alpha'' from
the strong, unsupported claim ``re-estimating a rolling ranking creates
useful dynamic forecasts.'' Conflating the two is the central error the
whole design is built to prevent.

\section{What would count as evidence, and what the study does not claim}\label{rq-evidence}

Each question is tied to a layer of evidence from Chapter~1, and each has a
clear failure mode. Q1 is answered only if a \emph{dependence-robust}
joint test rejects $H_0:\alpha=0$ and individual alphas survive
false-discovery control; a rejection under an IID-Gaussian test alone
(GRS) would not suffice, because the diagnostics reject those assumptions
(Chapter~13). Q2 is answered by a rank \emph{distribution}; a point ranking
with no interval is not an answer. Q3 is answered only by a strictly
leakage-free forward test \emph{and} an incremental decomposition;
in-sample fit, or a forward correlation without the incremental
comparison, would leave $H_{\text{stable}}$ and $H_{\text{dynamic}}$
unseparated.

Finally, the study is explicit about its non-goals. It does not claim an
implementable trading profit --- the reported spreads are gross of costs and
short hard-to-borrow small-growth stocks; it does not claim causal skill;
and it does not claim that its local data exactly match the current
official source, a discrepancy left visible as \textsc{not\_comparable}
(Chapter~4). Stating what is \emph{not} concluded is part of answering the
research question honestly.

